\documentclass[11pt,a4paper]{article}
\usepackage[T1]{fontenc}
\usepackage{lmodern,amsmath,amssymb,booktabs,microtype}
\usepackage[margin=24mm]{geometry}
\usepackage[colorlinks=true,urlcolor=blue,linkcolor=blue]{hyperref}
\newcommand{\EE}{\mathbb E}
\newcommand{\PP}{\mathbb P}
\newcommand{\core}{\mathcal E}
\DeclareMathOperator{\tr}{tr}
\title{Rank filters in the quadratic memory core\\Positive product channels and double-scaled tail control}
\author{Alan Day \and Chengheng\thanks{AI research assistant (OpenAI).}}
\date{Local research note --- 6 November 2026}
\begin{document}\maketitle
\begin{abstract}
The previously identified finite dynamical residual is resolved into
input and output irreducible ranks. In the tested system, all eleven
nonzero blocks of the four quadratic maps are collinear: the residual
reflects different channel weights, rather than another multiplication
direction within a block. Independently, covariance composition gives
a positive output-rank law for products of isotropic rank channels.
Three exact Casimir moments imply double-scaled concentration and a
tail bound that can control filters of at most linear deficit growth.
The full memory filter and its spectral limit remain to be evaluated.
\end{abstract}

\section{Rank resolution without discarding multiplicity}
Use the standard-CAR, overlapping complex Gaussian $p$-body ensemble,
with $\EE J_{IJ}J_{KL}=w\delta_{IL}\delta_{JK}$ and $\hbar=1$.
At charge $m$, the scalar product is $\EE\tr_m X^\dagger Y$,
where $\tr_m$ is normalized. The target remains
\begin{equation}
 \core=\sum_i(\operatorname{ad}_H^2c_i)^\dagger
                    \operatorname{ad}_H^2c_i
              +\tfrac14[H,[H,A]],\quad
 A=\sum_i[H,c_i]^\dagger[H,c_i].
\end{equation}
The preceding exact contraction gives $Q_\core=\Delta\core/2$ and
$W_2\core=Q_\core-2\EE\core$. Here $W_2$ is coupling chaos.

Let $E_{ij}=c_i^\dagger c_j$. On neutral operators the Casimir is
\begin{equation}
 \mathcal C X=\sum_{ij}[E_{ij},[E_{ji},X]],\qquad
 c_r=2r(N+1-r),\quad 0\le r\le R:=\min(m,N-m).
\end{equation}
The exact orthogonal projector is
\begin{equation}
 P_r=\prod_{s\ne r}\frac{\mathcal C-c_s}{c_r-c_s},\qquad
 d_r=\binom Nr^2-\binom N{r-1}^2.
\end{equation}
Use the analogous projectors in the $p$-particle coupling space on
both indices of the symmetric quadratic coefficient tensor.
Together with $P_r$ on the output they give orthogonal blocks
$(a,b;r)$, with $a\le b$. The sum of their positive Gram matrices
is the original Gram. Tensor-product multiplicities are retained
inside each block. This is essential in the established reduced
product formula of Kota--Vyas \cite{KV}, Eq. (35): a triple of
rank labels alone does not generally specify an intertwiner.

\newpage
\section{What the finite residual actually contains}
For $N=5,p=2,m=3,w=1$, resolve the four maps
\begin{equation}
 B=H^2-CI,\quad O=A-a_0I,\quad J_HF=\{H,F\}/2,
 \quad W_2\core.
\end{equation}
Here $C=18$, $a_0=\EE\tr_m A$, and $F=W_1T_3+86H$ is the
previous linear rank-filtered field. All 18 blocks are evaluated
without choosing Clebsch--Gordan bases. Seven vanish. In every
nonzero block the Gram of these four maps has rank one.

Define $H_a$ by projecting the same coupling realization onto
rank $a$. Put
$Z_{aa;r}=P_r(H_a^2)$ and
$Z_{ab;r}=P_r(\{H_a,H_b\})$ for $a<b$.
Then the raw maps have the exact realization-wise reconstruction
\begin{equation}
 (H^2,A,J_HF,Q_\core)
      =\sum_{a\le b,r}(1,o_{ab;r},f_{ab;r},g_{ab;r})Z_{ab;r}.
 \label{eq:filters}
\end{equation}
The complete table is
\begin{center}
\begin{tabular}{ccrrrr}
\toprule
$(a,b)$ & $r$ & $\|W_2Z_{ab;r}\|^2$ & $o_{ab;r}$ & $f_{ab;r}$ & $g_{ab;r}$\\
\midrule
$(0,0)$&0&$81/50$&$4/3$&$-10$&$128$\\
$(0,1)$&1&$864/25$&$4/3$&$-15/4$&$334/3$\\
$(0,2)$&2&$27$&$4/3$&$-6$&$352/3$\\
$(1,1)$&0&$192/25$&$7/4$&$5/2$&$148$\\
$(1,1)$&1&$448/75$&$1/2$&$5/2$&$88$\\
$(1,1)$&2&$320/3$&$5/4$&$5/2$&$409/4$\\
$(1,2)$&1&$160/3$&$2$&$1/4$&$595/4$\\
$(1,2)$&2&$224/3$&$1/2$&$1/4$&$647/8$\\
$(2,2)$&0&$3/2$&$4$&$-2$&$400$\\
$(2,2)$&1&$28/3$&$-1$&$-2$&$-20$\\
$(2,2)$&2&$140/3$&$2$&$-2$&$154$\\
\bottomrule
\end{tabular}
\end{center}
Subtracting the old three global directions leaves a positive norm
because $g$ is not one global linear combination of $1,o,f$ across
all $r>0$ blocks. Its squared norm remains
$248945355/14903$, agreeing exactly with the preceding calculation.
Output rank zero is precisely the sample-dependent scalar part.

The rank-one assertion follows from the full coefficient Gram,
not a fit to a few matrices. Two separate integer realizations
also verify all eight operator identities in \eqref{eq:filters}.
This result does not make the representation tensor products
multiplicity-free, and does not establish collinearity at general
$N,p,m$. It identifies the source of the known finite witness.

\newpage
\section{A positive law for the output rank of a product}
Now consider independent isotropic Hermitian rank-$a$ and rank-$b$
operators $U,V$ on $\mathcal H_m$, normalized by
$\EE U^2=\EE V^2=I$. Write $d=\binom Nm$ and define
\begin{equation}
 \pi_{ab}(r)=\EE\|P_r(UV)\|_{\tr_m}^2.
\end{equation}
These are nonnegative, sum to one, and vanish for $r>a+b$.
This auxiliary product is used to control ranks; it is not being
substituted for the same-coupling memory operator.

Let $\Phi_a(X)=\EE UXU$. Its rank-$s$ eigenvalue is the known
Racah crossing factor, in the conventions of our earlier covariance
calculation and the terminating definition \cite{DLMF},
\begin{equation}
 \chi_a(s)={}_4F_3\!\left(\begin{matrix}-a,a-N-1,-s,s-N-1\\
                                  -m,m-N,-N-1\end{matrix};1\right).
 \label{eq:chi}
\end{equation}
Since $\EE (UV)^\dagger X(UV)=\Phi_b\Phi_a(X)$, isotropy gives
\begin{equation}
 \chi_a(s)\chi_b(s)=\sum_r\pi_{ab}(r)\chi_r(s).
 \label{eq:composition}
\end{equation}
Orthogonality of the covariance maps is equivalently
\begin{equation}
 \sum_{s=0}^R d_s\chi_a(s)\chi_b(s)=\frac{d^2}{d_a}\delta_{ab}.
\end{equation}
It follows either from the standard Racah transform or by reshuffling
the orthogonal rank projectors in their Choi matrices. Therefore
\begin{equation}
 \boxed{\pi_{ab}(r)=\frac{d_r}{d^2}
       \sum_{s=0}^R d_s\chi_a(s)\chi_b(s)\chi_r(s).}
 \label{eq:law}
\end{equation}
The positivity in \eqref{eq:law} follows from the matrix-product
construction even though the displayed summands need not be positive.
This is the standard covariance/Racah machinery applied to the
output-rank distribution; it does not require individual multiplicity
amplitudes or their phases.

For $N=5,m=3$, independent full CAR calculations give
\begin{equation}
 \pi_{12}=(0,5/27,22/27),\qquad
 \pi_{22}=(1/75,176/675,98/135).
\end{equation}
They agree with \eqref{eq:law} using all 10,000 ordered coordinate
pairs in each case and every output rank.

\newpage
\section{Exact moments and the double-scaled deficit}
Set $x_r=r(N+1-r)$ and $C_f=m(N-m)(N+1)/N$. The $s=1$ case of
\eqref{eq:composition} gives the exact identity
\begin{equation}
 \EE x_r=x_a+x_b-\frac{x_ax_b}{C_f}.
\end{equation}
For higher moments define the monic polynomials and coefficients
\begin{align}
 L_j(x)&=\prod_{\ell=0}^{j-1}[x-\ell(N+1-\ell)],\\
 A_{nj}&=\frac{(-n)_j(n-N-1)_j}
 {(-m)_j(m-N)_j(-N-1)_j\,j!},\qquad
 \chi_r(n)=\sum_{j=0}^n A_{nj}L_j(x_r).
\end{align}
Writing $M_j=\EE L_j(x_r)$ gives a triangular exact recursion
\begin{equation}
 M_0=1,\qquad
 M_n=\frac{\chi_a(n)\chi_b(n)-\sum_{j<n}A_{nj}M_j}{A_{nn}},
 \quad n\le R. \label{eq:recursion}
\end{equation}
No unknown recoupling amplitude occurs in these summed moments.

Take $N=p^2/\lambda$, $m=\zeta N$,
$a=\alpha p$, $b=\beta p$, along admissible integer sequences.
Here $\lambda>0$ is fixed, $0\le\alpha,\beta\le1$, and $\zeta$
stays in a compact subset of $(0,1)$. Put
\begin{equation}
 t=a+b,\quad S=t-r,\quad D_x=x_t-x_r,
 \quad \theta=\alpha\beta\lambda
            \left[\frac1{\zeta(1-\zeta)}-2\right].
\end{equation}
Evaluation of \eqref{eq:recursion} through order three yields
\begin{equation}
 \left(\frac{\EE D_x}{N},\frac{\EE D_x^2}{N^2},
                    \frac{\EE D_x^3}{N^3}\right)
 \longrightarrow(\theta,\theta^2+\theta,
                         \theta^3+3\theta^2+\theta).
 \label{eq:three}
\end{equation}
This follows by exact rational algebra, retaining symbolic
$\alpha,\beta,\zeta,\lambda$. The denominators depend only on
$N,m$; after scaling they are bounded away from zero on the stated
compact sets. The rational degree comparison therefore also gives
uniform convergence in the two input-rank fractions and filling.

Moreover $0\le S\le t\le2p$ and
\begin{equation}
 D_x=S(N+1-2t+S).
\end{equation}
For $S>0$, $D_x/(NS)=1+O(p/N)$ uniformly. Thus \eqref{eq:three}
also gives the first three moments of $S$. In particular, with fixed
input deficits $a=p-u,b=p-v$ and half filling,
\begin{equation}
 \EE S\to2\lambda,\qquad
 \operatorname{Var}S\to2\lambda,\qquad
 \EE S^3\to8\lambda^3+12\lambda^2+2\lambda.
\end{equation}
These three moments agree with a Poisson law, but do not prove that law.

\newpage
\section{Tail control for the same-coupling quadratic problem}
The preceding moments imply the useful uniform estimates
\begin{align}
 \limsup\PP(S\ge K)&\le\frac{\theta^2+\theta}{K^2},\\
 \limsup\EE\big[(1+S)^2\mathbf1_{S\ge K}\big]
 &\le\frac{4(\theta^3+3\theta^2+\theta)}K,
 \qquad K\ge1.
 \label{eq:tails}
\end{align}
The second line uses $(1+S)^2\le4S^2\le4S^3/K$ on the indicated
tail. It controls squared filters with at most linear growth in $S$.

To connect this independent-product bound to the actual Gaussian
quadratic tensor, decompose $H=\sum_a H_a$ into its independent
irreducible Gaussian components, with $\EE H_a^2=C_aI$.
For $a<b$, the triangle inequality in each output-rank subspace gives
\begin{equation}
 \EE\|P_r\{H_a,H_b\}\|^2
       \le4C_aC_b\pi_{ab}(r).
\end{equation}
For $a=b$, Gaussian quadratic decoupling introduces an independent
copy $H_a'$ and gives
\begin{align}
 \EE\|P_rW_2(H_a^2)\|^2
 &=\tfrac12\EE\|P_r\{H_a,H_a'\}\|^2\\
 &\le2C_a^2\pi_{aa}(r).
\end{align}
These inequalities retain the exchange interference rather than
equating an ordered product with a symmetrized same-coupling field.
The known input weights $C_a/C=\omega_a$ have factorial tails in
the rank deficit $p-a$. Combining that input bound with
\eqref{eq:tails} shows tightness of the output deficit $2p-r$ for
the unfiltered $W_2(H^2)$ coefficient measure: first truncate the
input deficits, then the conditional output deficit.

For the memory problem this provides a concrete conditional route.
If its reduced filter amplitudes on bounded input-deficit windows
are uniformly at most linear in $S$, then \eqref{eq:tails} controls
their squared output tails by $O(K^{-1})$. The exact finite table
does not yet prove this growth estimate at general $N,p,m$.
Nor does the independent-product law determine exchange phases
or multiplicity-resolved amplitudes of the full memory core.

The next calculation is therefore the general reduced filter,
including multiplicity where present, followed by its weighted
norm limit. Only after that step can the fourth-chaos descendants
be restored and combined into the genuine $Z_2$ spectrum. The
two finite-memory-parameter covariance, spatial reconstruction,
and the bulk meson mass dictionary remain subsequent objectives.

\newpage
\section{Evidence, conventions and retained corrections}
All numerical records are local and use the same physical conventions
as the preceding contraction note. The finite tensor calculation uses
$w=1$; the product probabilities are variance-normalized. Neither is
a direct assignment of de Sitter masses or finite-$G$ dynamics.

The exact rank computation checks all projectors for idempotence,
orthogonality, eigenvalue and dimension. Its 18 block Grams sum
exactly to both the preceding full Gram and, after removing $r=0$,
the preceding sample-traceless Gram. The largest saved integer
intermediate is $710400$. Gram summation uses bounded integer
chunks and rational denominators. The full coefficient tensors,
not disorder samples, establish the eleven collinearity identities.

The output law is checked by four finite Racah transforms and
their orthogonality relations. The factorial recursion is checked
only at admissible probe orders. Two independent projected-CAR
product calculations verify all output probabilities and the
symmetrization bound. Eight large-$p$ cases evaluate the exact
moment formulas at 80 and 120 decimal digits; saved 65-digit
upper bounds agree to the specified $10^{-60}$ relative threshold.
These precision checks are not interval rounding certificates.

For example, at $\lambda=1/4$, $a=b=p$, the exact-arithmetic upper
bound $\EE S^2\le\EE D_x^2/(N+1-4p)^2$ is approximately
$0.73772646$ at $p=16$ and $0.74809847$ at $p=128$; its limit
is $3/4$. The actual finite deviations are retained.

Three verification corrections are preserved with their original
sources. A list was first used as a NumPy advanced index. A symbolic
reconstruction initially used structural equality before expanding
its exact zero residuals. A moment check initially requested probe
order three when $R=2$; its admissible domain was corrected, while
the larger systems retain the full third-order check. No saved
scientific result or tolerance was changed to obtain agreement.

\begin{thebibliography}{9}
\bibitem{KV} V. K. B. Kota and M. Vyas,
\emph{Random Matrix Theory for Transition Strength Densities in Finite
Quantum Systems: Results from Embedded Unitary Ensembles},
\href{https://arxiv.org/pdf/1504.00759}{arXiv:1504.00759},
Sec. 3, Eqs. (34)--(35). The reduced product retains a multiplicity label.
\bibitem{Kota} V. K. B. Kota,
\emph{Bivariate moments of the two-point correlation function for embedded
Gaussian unitary ensemble with $k$-body interactions},
\href{https://arxiv.org/html/2208.11312}{arXiv:2208.11312v2},
Sec. IV, covariance decomposition.
\bibitem{DLMF} NIST Digital Library of Mathematical Functions,
\href{https://dlmf.nist.gov/18.26.E3}{Eq. 18.26.3},
terminating Racah-polynomial definition. The parameter identification
is the one used in the project's earlier rank-covariance calculation.
\end{thebibliography}
\end{document}
